\documentclass[../../main.tex]{subfiles}
\begin{document}
\section{基本定义}\label{section:ch07-基本定义}

一个\textbf{Lie 群}是一个光滑流形 $G$（无边的），它同时也是代数意义下的群，并且乘法映射 $m\colon G\times G\to G$ 与取逆映射 $i\colon G\to G$，即
\begin{align*}
m(g,h)=gh,\qquad i(g)=g^{-1},
\end{align*}
都是光滑的。Lie 群特别地是一个\textbf{拓扑群}（即一个具有群结构的拓扑空间，使得乘法映射与取逆映射都连续）。

任意 Lie 群中的群运算以并置记号表示，但在某些交换群中，例如 $\mathbb{R}^n$，运算通常写成加法形式。传统上用符号 $e$ 表示任意 Lie 群的单位元（源自德语 \emph{Einselement}，"单位元"），我们遵循这一约定，但在有更常用记号的具体例子中例外（例如矩阵群中的单位矩阵记作 $I_n$，或 $\mathbb{R}^n$ 中的单位元记作 $0$）。

光滑性条件的如下另一种刻画有时很有用。（更强的结果另见\cref{exercise:ch07-03}。）

\begin{proposition}\label{proposition:ch07-01}
若 $G$ 是一个具有群结构的光滑流形，且由 $(g,h)\mapsto gh^{-1}$ 给出的映射 $G\times G\to G$ 是光滑的，则 $G$ 是一个 Lie 群。
\end{proposition}

\begin{exercise}\label{exercise:ch07-02}
证明\cref{proposition:ch07-01}。
\end{exercise}

若 $G$ 是一个 Lie 群，则任意元素 $g\in G$ 定义了映射 $L_g,R_g\colon G\to G$，分别称为\textbf{左平移}与\textbf{右平移}，即
\begin{align*}
L_g(h)=gh,\qquad R_g(h)=hg.
\end{align*}
由于 $L_g$ 可以表示为光滑映射的复合
\begin{align*}
G\xrightarrow{\iota_g}G\times G\xrightarrow{m}G,
\end{align*}
其中 $\iota_g(h)=(g,h)$，$m$ 为乘法映射，由此可知 $L_g$ 是光滑的。它实际上是 $G$ 的微分同胚，因为 $L_{g^{-1}}$ 是其光滑逆。类似地，$R_g\colon G\to G$ 是微分同胚。正如我们在下文中将反复看到的，Lie 群的许多重要性质都源于这样一个事实：我们可以通过这样的整体微分同胚系统地将任意点映射到任意其他点。

\begin{example}\label{example:ch07-03}[Lie 群]
下列每个流形在指定的群运算下都是 Lie 群。
\begin{enumerate}[(a)]
\item 一般线性群 $\operatorname{GL}(n,\mathbb{R})$ 是行列式非零的 $n\times n$ 实矩阵的集合。它在矩阵乘法下构成群，并且是向量空间 $\mathcal{M}(n,\mathbb{R})$ 的开子流形，如\cref{example:ch01-27}中所述。乘法是光滑的，因为乘积矩阵 $AB$ 的矩阵元是 $A$ 和 $B$ 的矩阵元的多项式。取逆由 Cramer 法则可知是光滑的。

\item 记 $\operatorname{GL}^{+}(n,\mathbb{R})$ 为 $\operatorname{GL}(n,\mathbb{R})$ 中行列式为正的矩阵组成的子集。由于 $\det(AB)=(\det A)(\det B)$ 且 $\det(A^{-1})=1/\det A$，它是 $\operatorname{GL}(n,\mathbb{R})$ 的子群；又由于它是连续行列式函数在 $(0,\infty)$ 下的原像，它是 $\operatorname{GL}(n,\mathbb{R})$ 的开子集，因而是一个 $n^2$ 维流形。群运算是 $\operatorname{GL}(n,\mathbb{R})$ 中群运算的限制，因此是光滑的。故 $\operatorname{GL}^{+}(n,\mathbb{R})$ 是一个 Lie 群。

\item 设 $G$ 是任意 Lie 群，$H\subseteq G$ 是一个开子群（即同时也是开子集的子群）。由与 (b) 部分相同的论证，$H$ 在继承的群结构与光滑流形结构下是一个 Lie 群。

\item 复一般线性群 $\operatorname{GL}(n,\mathbb{C})$ 是在矩阵乘法下可逆的 $n\times n$ 复矩阵构成的群。它是 $\mathcal{M}(n,\mathbb{C})$ 的开子流形，因而是一个 $2n^2$ 维光滑流形；它是 Lie 群，因为矩阵乘积与逆是矩阵元实部和虚部的光滑函数。

\item 若 $V$ 是任意实或复向量空间，$\operatorname{GL}(V)$ 表示从 $V$ 到自身的可逆线性映射的集合。它在复合下构成群。若 $V$ 为有限维 $n$ 维，则 $V$ 的任意一组基决定了 $\operatorname{GL}(V)$ 与 $\operatorname{GL}(n,\mathbb{R})$ 或 $\operatorname{GL}(n,\mathbb{C})$ 之间的同构，故 $\operatorname{GL}(V)$ 是一个 Lie 群。任意两个这样的同构之间的转移映射由形如 $A\mapsto BAB^{-1}$ 的映射给出（其中 $B$ 是两组基之间的转移矩阵），该映射是光滑的。因此，$\operatorname{GL}(V)$ 上的光滑流形结构与基的选取无关。

\item 实数域 $\mathbb{R}$ 与 Euclid 空间 $\mathbb{R}^n$ 在加法下是 Lie 群，因为 $x+y$ 的坐标是 $(x,y)$ 的光滑（线性！）函数。

\item 类似地，$\mathbb{C}$ 与 $\mathbb{C}^n$ 在加法下是 Lie 群。

\item 非零实数集合 $\mathbb{R}^*$ 在乘法下是一个一维 Lie 群。（事实上，若将 $1\times 1$ 矩阵与相应实数等同，则它恰好是 $\operatorname{GL}(1,\mathbb{R})$。）正实数子集 $\mathbb{R}^+$ 是开子群，因而是一个一维 Lie 群。

\item 非零复数集合 $\mathbb{C}^*$ 在复数乘法下是一个二维 Lie 群，可与 $\operatorname{GL}(1,\mathbb{C})$ 等同。

\item 圆周 $S^1\subseteq\mathbb{C}^*$ 在复数乘法下是光滑流形和群。以适当的角函数作为 $S^1$ 开子集上的局部坐标（见\cref{exercise:ch01-08}），乘法与取逆具有光滑坐标表达式 $(\theta_1,\theta_2)\mapsto\theta_1+\theta_2$ 与 $\theta\mapsto-\theta$，因此 $S^1$ 是一个 Lie 群，称为\textbf{圆周群}。

\item 给定 Lie 群 $G_1,\dots,G_k$，它们的直积是乘积流形 $G_1\times\cdots\times G_k$，其群结构由按分量乘法给出：
\begin{align*}
(g_1,\dots,g_k)(g_1',\dots,g_k')=(g_1g_1',\dots,g_kg_k').
\end{align*}
容易验证它是 Lie 群。

\item $n$ 维环面 $\mathbb{T}^n=S^1\times\cdots\times S^1$ 是一个 $n$ 维交换 Lie 群。

\item 任意赋予离散拓扑的群都是拓扑群，称为\textbf{离散群}。若该群还是有限或可数无限的，则它是一个零维 Lie 群，称为\textbf{离散 Lie 群}。
\end{enumerate}
\end{example}

\end{document}